\section{Phát hiện outlier IQR [Nâng cao]}
\subsection{Phân tích \& Thiết kế (M0)}

\begin{itemize}

    \item \textbf{Input:}

    File \texttt{values.csv} chứa:

    \begin{itemize}
        \item Mỗi dòng là một số thực
        \item Không có missing
    \end{itemize}

    \item \textbf{Xử lý:}

    \begin{itemize}

        \item Sắp xếp dữ liệu tăng dần

        \item Tính:
        \[
            Q1,\ Q3
        \]

        \item Công thức:
        \[
            IQR = Q3 - Q1
        \]

        \item Outlier nếu:
        \[
            x < Q1 - 1.5 \times IQR
        \]

        hoặc

        \[
            x > Q3 + 1.5 \times IQR
        \]

    \end{itemize}

    \begin{itemize}

        \item \textbf{Case biên:}

        \begin{itemize}
            \item File rỗng
            \[
                \Rightarrow
                Q1=Q3=IQR=0
            \]

            \item Tất cả giá trị giống nhau
            \[
                \Rightarrow IQR=0
            \]

            \item Không có outlier
            \[
                \Rightarrow
                outliers.csv \text{ chỉ có header}
            \]
        \end{itemize}

        \item \textbf{Module 1 (Đọc file):}

        \begin{itemize}
            \item Hàm \texttt{readValues()}
        \end{itemize}

        \item \textbf{Module 2 (Tính percentile):}

        \begin{itemize}
            \item Hàm \texttt{computeQuartile()}
            \item Tính:
            \[
                Q1,\ Q3
            \]
        \end{itemize}

        \item \textbf{Module 3 (Phát hiện outlier):}

        \begin{itemize}
            \item Hàm \texttt{detectOutliers()}
        \end{itemize}

        \item \textbf{Module 4 (Ghi outlier):}

        \begin{itemize}
            \item Hàm \texttt{writeOutliers()}
        \end{itemize}

        \item \textbf{Module 5 (Ghi summary):}

        \begin{itemize}
            \item Hàm \texttt{writeSummary()}
        \end{itemize}

    \end{itemize}

    \item \textbf{Output:}

    File \texttt{outliers.csv}

\begin{lstlisting}
index,value
...
\end{lstlisting}

    File \texttt{summary.txt}

\begin{lstlisting}
Q1=<f>
Q3=<f>
IQR=<f>
n_outliers=<n>
\end{lstlisting}

    Dùng:
\begin{lstlisting}
setprecision(6)
\end{lstlisting}

\end{itemize}

\subsection{Mã nguồn C++11 (Tách module)}

\begin{lstlisting}
#include <iostream>
#include <fstream>
#include <vector>
#include <algorithm>
#include <iomanip>

using namespace std;

bool readValues(
    const string& path,
    vector<double>& values
) {
    ifstream fin(path);

    if (!fin) return false;

    double x;

    while (fin >> x) {
        values.push_back(x);
    }

    return true;
}

double computeQuartile(
    vector<double> sorted,
    double q
) {
    if (sorted.empty()) return 0;

    double pos =
        q * (sorted.size() - 1);

    int left = (int)pos;

    int right = left + 1;

    if (right >= (int)sorted.size()) {
        return sorted[left];
    }

    double frac = pos - left;

    return sorted[left]
           +
           frac
           *
           (sorted[right] - sorted[left]);
}

vector<pair<int, double>> detectOutliers(
    const vector<double>& values,
    double Q1,
    double Q3
) {
    vector<pair<int, double>> outliers;

    double IQR = Q3 - Q1;

    double lower =
        Q1 - 1.5 * IQR;

    double upper =
        Q3 + 1.5 * IQR;

    for (size_t i = 0; i < values.size(); ++i) {

        if (values[i] < lower
            || values[i] > upper) {

            outliers.push_back(
                {i, values[i]}
            );
        }
    }

    return outliers;
}

bool writeOutliers(
    const string& path,
    const vector<pair<int, double>>& outliers
) {
    ofstream fout(path);

    if (!fout) return false;

    fout << fixed << setprecision(6);

    fout << "index,value\n";

    for (size_t i = 0; i < outliers.size(); ++i) {

        fout << outliers[i].first
             << ","
             << outliers[i].second
             << "\n";
    }

    return true;
}

bool writeSummary(
    const string& path,
    double Q1,
    double Q3,
    double IQR,
    int nOutliers
) {
    ofstream fout(path);

    if (!fout) return false;

    fout << fixed << setprecision(6);

    fout << "Q1=" << Q1 << "\n";

    fout << "Q3=" << Q3 << "\n";

    fout << "IQR=" << IQR << "\n";

    fout << "n_outliers="
         << nOutliers
         << "\n";

    return true;
}

int main() {

    vector<double> values;

    if (!readValues(
        "values.csv",
        values
    )) {
        cerr << "Loi mo file input!\n";
        return 1;
    }

    vector<double> sorted =
        values;

    sort(
        sorted.begin(),
        sorted.end()
    );

    double Q1 =
        computeQuartile(sorted, 0.25);

    double Q3 =
        computeQuartile(sorted, 0.75);

    double IQR =
        Q3 - Q1;

    vector<pair<int, double>> outliers =
        detectOutliers(
            values,
            Q1,
            Q3
        );

    if (!writeOutliers(
        "outliers.csv",
        outliers
    )) {
        cerr << "Loi mo file outlier!\n";
        return 1;
    }

    if (!writeSummary(
        "summary.txt",
        Q1,
        Q3,
        IQR,
        outliers.size()
    )) {
        cerr << "Loi mo file summary!\n";
        return 1;
    }

    return 0;
}
\end{lstlisting}

\subsection{Kế hoạch kiểm thử (Test Cases)}

\textbf{Test Case 1 (Có outlier)}

\textbf{Input (values.csv):}

\begin{lstlisting}
10
11
12
13
14
100
\end{lstlisting}

\textbf{Expected Output (outliers.csv):}

\begin{lstlisting}
index,value
5,100.000000
\end{lstlisting}

\textbf{Expected Output (summary.txt):}

\begin{lstlisting}
Q1=11.250000
Q3=13.750000
IQR=2.500000
n_outliers=1
\end{lstlisting}

\vspace{0.5cm}

\textbf{Test Case 2 (Không có outlier)}

\textbf{Input (values.csv):}

\begin{lstlisting}
1
2
3
4
5
\end{lstlisting}

\textbf{Expected Output (outliers.csv):}

\begin{lstlisting}
index,value
\end{lstlisting}

\textbf{Expected Output (summary.txt):}

\begin{lstlisting}
Q1=2.000000
Q3=4.000000
IQR=2.000000
n_outliers=0
\end{lstlisting}

\vspace{0.5cm}

\textbf{Test Case 3 (File rỗng)}

\textbf{Input (values.csv):}

\begin{lstlisting}

\end{lstlisting}

\textbf{Expected Output (outliers.csv):}

\begin{lstlisting}
index,value
\end{lstlisting}

\textbf{Expected Output (summary.txt):}

\begin{lstlisting}
Q1=0.000000
Q3=0.000000
IQR=0.000000
n_outliers=0
\end{lstlisting}